% Part: proof-theory
% Chapter: normalization
% Section: reductions

\documentclass[../../../include/open-logic-section]{subfiles}

\begin{document}

\olfileid[pt]{pt}{nor}{red}

\olsection{Conversões de redução}

Se um corte em !!a{proof} em \Log{N1i} tiver comprimento~$1$, ele será
um segmento constituído por uma única ocorrência de !!{formula}, e essa
!!{formula} será tanto a conclusão de uma inferência de
!!a{introduction} quanto a premissa maior de uma inferência de
!!a{elimination}. Na prova da normalização, esses cortes são
``reduzidos'' substituindo-se a sub-!!{proof} que termina nesse par
!!{introduction}/!!{elimination} por uma nova !!{proof} na qual o par
de inferências foi eliminado. Essa substituição pode produzir novos
cortes. Contudo, podemos nos convencer de que os novos cortes têm posto
menor que o corte reduzido. Assim, a nova !!{proof} ou tem comprimento
de corte menor (porque um corte maximal de comprimento~$1$ foi removido)
ou tem o posto de corte da !!{proof} reduzido (se o corte maximal em questão era o
único corte maximal).

Por exemplo, se as inferências envolvidas forem \Intro{\lif} e
\Elim{\lif}, $!A \ident !B \lif !C$, e houver uma
sub-!!{proof}~$\delta_1$ que termine em:
\[
\AxiomC{$\Discharge{!B}{x}$}
\RightLabel{$\delta_2$}
\DeduceC{$!C$}
\DischargeRule{\Intro{\lif}}{x}
\UnaryInfC{$!B \lif !C$}
\AxiomC{}
\RightLabel{$\delta_3$}
\DeduceC{$!B$}
\RightLabel{\Elim{\lif}}
\BinaryInfC{$!C$}
\DisplayProof
\]
Para remover esse corte, substituímos a subprova~$\delta_1$ em~$\delta$
pelo seguinte:
\[
\AxiomC{}
\RightLabel{$\delta_3$}
\DeduceC{$!B$}
\RightLabel{$\Subst{\delta_2}{\delta_3}{!B^x}$}
\DeduceC{$!C$}
\DisplayProof
\]
Por $\Subst{\delta_2}{\delta_3}{!B^x}$ entendemos a !!{proof} que
resulta de~$\delta_2$ quando substituímos cada premissa~$!B$ rotulada
por~$x$ pela !!{proof} inteira~$\delta_3$. Essa !!{proof} já não tem
premissas abertas da forma~$!B$ com o rótulo~$x$. Ela não tem outras
premissas rotuladas por~$x$, pois supomos, como de costume, que duas
!!{formula}s distintas não ocorrem como premissas com o mesmo rótulo.
Portanto, a nova !!{proof} tem as mesmas premissas abertas
que~$\delta_1$; cada inferência em~$\delta_2$ que descarrega premissas
descarrega as mesmas premissas em~$\Subst{\delta_2}{\delta_3}{!B^x}$; e
cada inferência em~$\delta_3$ que descarrega premissas dá origem a
possivelmente várias cópias, que descarregam as premissas
correspondentes em~$\Subst{\delta_2}{\delta_3}{!B^x}$.

Todo corte que esteja inteiramente dentro de $\delta_2$, $\delta_3$ ou
$\delta_4$, ou inteiramente fora de~$\delta_1$ em~$\delta$, permanece
inalterado: ele ainda envolve as mesmas !!{formula}s e passa pelas
mesmas inferências; portanto, tem, em particular, o mesmo posto e o
mesmo comprimento que o corte correspondente em~$\delta$. Removemos um
corte maximal de comprimento~$1$ e, assim, reduzimos o comprimento de
corte ou reduzimos o posto de corte (se o comprimento de corte
de~$\delta$ era~$1$ e esse era o único corte maximal).

Contudo, duas coisas podem acontecer:
\begin{enumerate}
  \item Ao substituir todas as premissas $!B^x$ por $\delta_3$,
  multiplicamos todos os cortes em~$\delta_3$. Se alguns deles eram
  cortes maximais, o comprimento de corte da nova !!{proof} pode ser
  maior que o da antiga !!{proof}. Precisamos garantir que isso não
  aconteça. Fazemos isso aplicando conversões de redução somente aos
  cortes maximais situados \emph{mais à direita}, isto é, a cortes para
  os quais nenhum corte maximal se encontra em um ramo que passa
  por~$\delta$, à direita do corte que está sendo reduzido. Se houvesse
  um corte em~$\delta_3$, ele estaria à direita do corte reduzido aqui;
  portanto, não estaríamos trabalhando com um corte situado mais à
  direita. Ao escolher sempre os cortes mais à direita, garantimos que
  reduziremos o comprimento de corte da !!{proof} ou até mesmo seu posto
  de corte.
  \item Se uma premissa $!B^x$ em $\delta_2$ for a premissa maior de
  !!a{elimination}, e a última inferência de~$\delta_3$ for
  \Elim{\lor}, \Elim{\lexists} ou !!a{introduction}, teremos produzido
  um novo corte enxertando~$\delta_3$ em $\delta_2$ nessa premissa. O
  que era apenas uma premissa em~$\delta_2$ agora é um corte de
  comprimento~$1$ (a conclusão de !!a{introduction} e a premissa maior
  de !!a{elimination}) ou a última ocorrência de !!{formula} de um
  segmento de corte (premissa maior de !!a{elimination} e conclusão
  de~\Elim{\lor} ou \Elim{\lexists}). Se $B^x$ era uma premissa menor
  de \Elim{\lor} ou \Elim{\lexists}, ela já era a primeira ocorrência
  de !!{formula} de um segmento, e possivelmente de um segmento de
  corte. Na nova !!{proof}, esse segmento pode agora ser mais longo.
  Contudo, observe que a !!{formula} que constitui esse novo corte ou o
  segmento de corte existente é~$!B$, e $\depth{!B} < \depth{!B \lif
  !C}$. Portanto, embora possamos ter acrescentado cortes ou aumentado
  o comprimento de cortes, nenhum deles é um corte maximal. Assim,
  reduzimos o posto de corte ou, se restarem cortes do mesmo posto,
  reduzimos o comprimento de corte.
\end{enumerate}

O processo de reduzir um corte de comprimento~$1$ é chamado de
\emph{conversão de redução} (ou \emph{conversão de desvio}). Os padrões
de alguns dos pares de regras são apresentados em
\olref{tab:reduction-conversions}. (Os restantes ficam como exercícios.)

\begin{table}
\begin{multline*}
\bottomAlignProof
\AxiomC{$\Discharge{!B}{x}$}
\RightLabel{$\delta_2$}
\shortDeduce
\DeduceC{$!C$}
\DischargeRule{\Intro{\lif}}{x}
\UnaryInfC{$!B \lif !C$}
\AxiomC{}
\RightLabel{$\delta_3$}
\shortDeduce
\DeduceC{$!B$}
\RightLabel{\Elim{\lif}}
\BinaryInfC{$!C$}
\DisplayProof
\quad \longrightarrow\quad
\shoveright{\bottomAlignProof
\AxiomC{}
\RightLabel{$\delta_3$}
\shortDeduce
\DeduceC{$!B$}
\RightLabel{$\Subst{\delta_2}{\delta_3}{!B^x}$}
\shortDeduce
\DeduceC{$!C$}
\DisplayProof}
\\[4ex]
\shoveleft{\bottomAlignProof
\AxiomC{}
\RightLabel{$\delta_2$}
\shortDeduce
\DeduceC{$!B$}
\AxiomC{}
\RightLabel{$\delta_3$}
\shortDeduce
\DeduceC{$!C$}
\RightLabel{\Intro{\land}}
\BinaryInfC{$!B \land !C$}
\RightLabel{\Elim{\land}[1]}
\UnaryInfC{$!B$}
\DisplayProof
\quad \longrightarrow\quad
\bottomAlignProof
\AxiomC{}
\RightLabel{$\delta_2$}
\shortDeduce
\DeduceC{$!B$}
\DisplayProof}\qquad
\shoveright{\bottomAlignProof
\AxiomC{}
\RightLabel{$\delta_2$}
\shortDeduce
\DeduceC{$!B$}
\AxiomC{}
\RightLabel{$\delta_3$}
\shortDeduce
\DeduceC{$!C$}
\RightLabel{\Intro{\land}}
\BinaryInfC{$!B \land !C$}
\RightLabel{\Elim{\land}[2]}
\UnaryInfC{$!C$}
\DisplayProof
\quad \longrightarrow\quad
\bottomAlignProof
\AxiomC{}
\RightLabel{$\delta_3$}
\shortDeduce
\DeduceC{$!C$}
\DisplayProof}
\\[4ex]
\bottomAlignProof
\AxiomC{}
\RightLabel{$\delta_2$}
\shortDeduce
\DeduceC{$!B$}
\RightLabel{\Intro{\lor}[1]}
\UnaryInfC{$!B \lor !C$}
\AxiomC{$\Discharge{!B}{x}$}
\RightLabel{$\delta_3$}
\shortDeduce
\DeduceC{$!D$}
\AxiomC{$\Discharge{!C}{y}$}
\RightLabel{$\delta_4$}
\shortDeduce
\DeduceC{$!D$}
\DischargeRule{\Elim{\lor}}{x\,y}
\TrinaryInfC{$!D$}
\DisplayProof
\quad \longrightarrow\quad
\shoveright{\bottomAlignProof
\AxiomC{}
\RightLabel{$\delta_2$}
\shortDeduce
\DeduceC{$!B$}
\RightLabel{$\Subst{\delta_3}{\delta_2}{!B^x}$}
\shortDeduce
\DeduceC{$!D$}
\DisplayProof}
\\[2ex]
\shoveleft{\bottomAlignProof
\AxiomC{}
\RightLabel{$\delta_2$}
\shortDeduce
\DeduceC{$!B(c)$}
\RightLabel{\Intro{\lforall}}
\UnaryInfC{$\lforall[x][!B(x)]$}
\RightLabel{\Elim{\lforall}}
\UnaryInfC{$!B(t)$}
\DisplayProof
\quad \longrightarrow\quad
\bottomAlignProof
\AxiomC{}
\RightLabel{$\Subst{\delta_2}{t}{c}$}
\shortDeduce
\DeduceC{$!B(t)$}
\DisplayProof}
\\[4ex]
\shoveright{\bottomAlignProof
\AxiomC{}
\RightLabel{$\delta_2$}
\shortDeduce
\DeduceC{$!B(t)$}
\RightLabel{\Intro{\lexists}}
\UnaryInfC{$\lexists[x][!B(x)]$}
\AxiomC{$\Discharge{!B(c)}{x}$}
\RightLabel{$\delta_3$}
\shortDeduce
\DeduceC{$!D$}
\DischargeRule{\Elim{\lexists}}{x}
\BinaryInfC{$!D$}
\DisplayProof
\quad \longrightarrow\quad
\bottomAlignProof
\AxiomC{}
\RightLabel{$\delta_2$}
\shortDeduce
\DeduceC{$!B(t)$}
\RightLabel{$\Subst{\Subst{\delta_3}{t}{c}}{\delta_2}{!B(t)^x}$}
\shortDeduce
\DeduceC{$!D$}
\DisplayProof}
\end{multline*}
\caption{Algumas conversões de redução para \Log{N1i}.}
\ollabel{tab:reduction-conversions}
\end{table}

\begin{prob}
Apresente a conversão de redução para \Intro{\lnot} seguida de
\Elim{\lnot}.
\end{prob}

Resta considerar um caso. A definição de corte inclui o caso em que o
comprimento do corte é~$1$ e $!A$ é tanto a conclusão de~\FalseInt{}
quanto a premissa maior de uma inferência de !!a{elimination}. Esse caso
é tratado de maneira semelhante ao caso em que $!A$ é a conclusão de
uma regra de !!a{introduction}. Por exemplo, substitua
\[
\bottomAlignProof
\AxiomC{}
\RightLabel{$\delta_2$}
\DeduceC{$\lfalse$}
\UnaryInfC{$!B \lif !C$}
\AxiomC{}
\RightLabel{$\delta_3$}
\DeduceC{$!B$}
\RightLabel{\Elim{\lif}}
\BinaryInfC{$!C$}
\DisplayProof
\quad\text{por}\quad
\bottomAlignProof
\AxiomC{}
\RightLabel{$\delta_2$}
\DeduceC{$\lfalse$}
\UnaryInfC{$!C$}
\DisplayProof
\]
Precisamos ter um pouco de cuidado nos casos em que $!A \ident (!B
\lor !C)$ ou $!A \ident \lexists[x][!B(x)]$. Por exemplo, se
substituíssemos
\[
\bottomAlignProof
\AxiomC{}
\RightLabel{$\delta_2$}
\DeduceC{$\lfalse$}
\RightLabel{\FalseInt}
\UnaryInfC{$!B \lor !C$}
\AxiomC{$\Discharge{!B}{x}$}
\RightLabel{$\delta_3$}
\DeduceC{$!D$}
\AxiomC{$\Discharge{!C}{x}$}
\RightLabel{$\delta_4$}
\DeduceC{$!D$}
\RightLabel{\Elim{\lor}}
\TrinaryInfC{$!D$}
\DisplayProof
\quad\text{por}\quad
\bottomAlignProof
\AxiomC{}
\RightLabel{$\delta_2$}
\DeduceC{$\lfalse$}
\RightLabel{\FalseInt}
\UnaryInfC{$!D$}
\DisplayProof
\]
poderíamos estar introduzindo um novo corte cuja !!{formula} de corte
é~$!D$, e não há garantia de que $\depth{!D} < \depth{!B \lor !C}$!{}
Portanto, devemos proceder aqui da mesma maneira que na redução
\Intro{\lor}/\Elim{\lor} e substituir por
\[
\bottomAlignProof
\AxiomC{}
\RightLabel{$\delta_2$}
\DeduceC{$\lfalse$}
\RightLabel{\FalseInt}
\UnaryInfC{$!B$}
\RightLabel{$\delta_3$}
\DeduceC{$!D$}
\DisplayProof
\quad\text{ou}\quad
\bottomAlignProof
\AxiomC{}
\RightLabel{$\delta_2$}
\DeduceC{$\lfalse$}
\RightLabel{\FalseInt}
\UnaryInfC{$!C$}
\RightLabel{$\delta_4$}
\DeduceC{$!D$}
\DisplayProof
\]

\end{document}
